\section{Lean formalization: definitions, main lemmas, and completion path}
\label{sec:lean}

The archive contains an integration-oriented Lean outline under \texttt{lean/}.  The runtime used to prepare
this packet has no Lean executable, so the files were not compiled.  They contain no \texttt{axiom},
\texttt{sorry}, or \texttt{admit}; unfinished mathematics is represented by proof-bearing structures and
commented target signatures.  This section explains how to turn the outline into a complete formalization
against the checked \texttt{ConciseIPRA} library.

\subsection{Imported interfaces}

The kernel imports
\begin{lstlisting}[language={}]
import ConciseIPRA.PR
import ConciseIPRA.NativeBiInterpretation
\end{lstlisting}
and works in \texttt{LeanProofs.ConciseIPRA}.  The required checked declarations are:
\begin{itemize}
  \item \texttt{PR}, \texttt{PTm}, and \texttt{PRAAx};
  \item \texttt{Prf.subst}, \texttt{Prf.weaken}, and \texttt{Prf.trans};
  \item \texttt{PRA.qfInd};
  \item the logical constructors \texttt{Prf.impI}, \texttt{Prf.impE},
        \texttt{Prf.allE}, and \texttt{Prf.exI};
  \item \texttt{pra\_iff\_native}.
\end{itemize}
No semantic interpretation of \(\PR\) functions is needed by the generic compiler.

\subsection{The specialization kernel}

The principal theorem is written directly with current repository types:
\begin{lstlisting}[language={}]
theorem instantiateInto
    {k n : Nat}
    {Delta : List (Fm PTm k)} {A : Fm PTm k}
    (local : Delta |-[PRAAx] A)
    (sigma : Fin k -> PTm n)
    {Theta : List (Fm PTm n)}
    (available : Delta.map (fun B => B.subst sigma) <= Theta) :
    Theta |-[PRAAx] A.subst sigma :=
  (local.subst sigma).weaken available
\end{lstlisting}
The ASCII listing uses \texttt{|-[...]} and \texttt{<=} as visual stand-ins for the repository notation
\(\vdash[\cdot]\) and list inclusion; the supplied source uses the actual notation.

A proof schedule is
\begin{lstlisting}[language={}]
structure ProofBuild (Gamma : List (Fm PTm n))
    (done : List (Fm PTm n)) where
  derive : forall A, A in done -> Gamma |-[PRAAx] A
\end{lstlisting}
The \texttt{step} method cuts a local derivation from \texttt{done ++ Gamma} against stored proofs and base
hypotheses.  The arity-normalized node is:
\begin{lstlisting}[language={}]
structure NodeObligation (Gamma done : List (Fm PTm n)) where
  arity : Nat
  localContext : List (Fm PTm arity)
  localGoal : Fm PTm arity
  args : Fin arity -> PTm n
  localProof : localContext |-[PRAAx] localGoal
  available : localContext.map (fun B => B.subst args)
              <= done ++ Gamma
\end{lstlisting}
Its compiler is compositionally tiny:
\begin{lstlisting}[language={}]
def NodeObligation.compile
    (s : ProofBuild Gamma done)
    (N : NodeObligation Gamma done) :
    ProofBuild Gamma (N.goal :: done) :=
  s.step (instantiateInto N.localProof N.args N.available)
\end{lstlisting}
This definition is the formal content of Theorem~\ref{thm:staged-graph}.  A generated application file
repeatedly invokes it; the final theorem uses \texttt{ProofBuild.get}.

\subsection{Loop compilation}

The supplied \texttt{LoopObligation} record stores:
\begin{itemize}
  \item a formula \texttt{invariant : Fm PTm (n+1)};
  \item a proof \texttt{invariantQF : invariant.QF};
  \item a fuel term;
  \item actual \(\Prf\,\PRAAx\) base, step, and finish proofs.
\end{itemize}
The compiler invokes \texttt{PRA.qfInd}, eliminates the universal invariant at the fuel term, and uses
simultaneous cut for the finish step.  The intended first build task is to compile this generic file before
formalizing any new domain algorithm.  Any API mismatch should be a local spelling issue, since the proof
pattern was copied from the current checked interfaces and the existing kernel finite-use draft.

\subsection{Boolean primitive library}

Create \texttt{StagedContractKernel/BoolPR.lean}.  Define concrete symbols for:
\[
 \operatorname{isZero},\quad \dotminus,\quad \operatorname{eq},\quad
 \operatorname{lt},\quad \wedge_B,\quad\vee_B,\quad\neg_B,
\]
using \texttt{PR.comp} and \texttt{PR.prec}.  For each, prove equations at arbitrary arguments by combining
\texttt{PRA.comp}, \texttt{PRA.precZero}, and \texttt{PRA.precSucc}.  Then prove range lemmas such as
\[
 b(x)=0\vee b(x)=1
\]
and reflection lemmas.  These proofs use quantifier-free induction over the recursion counter and the
successor axioms.

It is useful to expose a common theorem shape:
\begin{lstlisting}[language={}]
structure ReflectedBool (n : Nat) where
  term : PTm n
  formula : Fm PTm n
  qf : formula.QF
  isBool : [] |-[PRAAx] (term = 0) or (term = 1)
  trueIff : [] |-[PRAAx] iff (term = 1) formula
\end{lstlisting}
The exact formula constructors must use the repository's \texttt{Fm} notation.  The structure is a
convenience; the final theorem remains a \texttt{Prf} term.

\subsection{Bounded-formula compiler}

The outline defines a host datatype \texttt{BForm n} with equality, order, propositional connectives, and
bounded quantifiers.  Add:
\begin{enumerate}
  \item \texttt{BForm.toFormula : BForm n -> Fm PTm n};
  \item \texttt{compileBForm : BForm n -> PR n};
  \item a recursive function returning a \texttt{BooleanizationResult}.
\end{enumerate}
For propositional constructors, compose the Boolean primitive reflection theorems.  For bounded quantifiers,
first prove generic finite-fold modules:
\begin{lstlisting}[language={}]
def allFold (p : PR (n+1)) : PR (n+1)
def anyFold (p : PR (n+1)) : PR (n+1)

theorem allFold_reflect ... : Prf PRAAx ...
theorem anyFold_reflect ... : Prf PRAAx ...
\end{lstlisting}
The fold invariant states the equivalence between the accumulator and the bounded quantifier below the
current counter.  Represent the bounded quantifier side by its already compiled Boolean, ensuring the
induction invariant itself is quantifier-free.

\subsection{Least search and enumeration}

Create \texttt{BoundedSearch.lean}.  The state may be a pair encoded by a fixed pairing \(\PR\) symbol or a
small vector code.  Define the transition from Section~\ref{sec:booleanization}.  The invariant should be
implemented as a conjunction of reflected Boolean equations:
\begin{align*}
 f=0 &\to \operatorname{allFailBelow}(t)=1,\\
 f=1 &\to w<t\wedge p(w)=1\wedge\operatorname{allFailBelow}(w)=1.
\end{align*}
Compile it through \texttt{LoopObligation}.  The final theorem should return a concrete witness term and
actual proofs of correctness and leastness.  Complete enumeration uses a bounded list append function and a
parallel exact-prefix invariant.

\subsection{Demand closure}

Create \texttt{Demand.lean} for each fixed vector arity, or write a generic list-coded version after the
fixed case compiles.  Define mixed-radix vectors, componentwise order, join, sum, and cap validity.  Given a
concrete transformer \(F\), the loop state is \((q,\mathit{stable})\).  The invariant includes:
\begin{enumerate}
  \item \(q\) is inside the cap;
  \item \(q^0\le q\);
  \item \(q\) is the current iterate;
  \item if not stable, the coordinate sum strictly increases at the next step.
\end{enumerate}
The fuel is the sum of remaining coordinate capacity plus one.  A separate local proof establishes
monotonicity of \(F\), enabling the least-fixed-point theorem.  Since \(F\) is an application-specific
concrete \(\PR\) term, its cap and monotonicity proofs are local node obligations, not generic axioms.

\subsection{Backward viability}

Create \texttt{Cloud.lean}.  For a fixed horizon, bounded candidate tables and edge verifiers are concrete
symbols.  Define a reverse table \(W\) by primitive recursion on distance from the final level.  Because the
\texttt{PR.prec} recursion argument increases, store the reversed level index \(h-i\).  Prove
Theorem~\ref{thm:viability} by induction on that distance.

The least path constructor is a forward trace.  Its invariant says that the current candidate is viable and
that the accumulated path prefix is valid.  The successor search uses \texttt{leastBoundedSearch\_correct}
on the conjunction of the edge verifier and next-level viability.  The defining equation of \(W\) provides
nonemptiness.

\subsection{Stable and diameter observables}

Exact branch erasure is already represented by \texttt{stableObservable}.  A richer module should define a
rational metric formula and prove:
\begin{lstlisting}[language={}]
theorem diameterSelected
  (selectedCorrect : ...)
  (diameter : ... |-[PRAAx] dist (obs c) (obs d) <= eps) :
  ... |-[PRAAx] dist (obs selected) (obs arbitrary) <= eps
\end{lstlisting}
Gap-based exactness is a finite triangle-inequality contradiction.  Provide reusable ordered-rational
lemmas so applications need only instantiate the labels, approximation radius, and separation bound.

\subsection{Resource library}

Create \texttt{MixedRadix.lean}.  For each fixed arity, define \texttt{pack}, projections, and
\texttt{rank}.  Prove the lower-coordinate geometric identity and lexicographic-rank lemma by induction on
the finite coordinate index at the host level, generating a theorem for the fixed arity.  A later generic
list-coded version may use a bounded fold and quantifier-free induction.

Height-to-list conversion should be a corollary of bounded enumeration.  Complete-list-to-height uses a
bounded maximum fold and the membership completeness proof.

\subsection{Generated application modules}

A domain compiler should generate a Lean module with this shape:
\begin{enumerate}
  \item concrete \texttt{PR} symbols and \texttt{PTm} abbreviations;
  \item local equations and correctness lemmas;
  \item a base context \(\Gamma\) exposing external boundary data;
  \item one \texttt{NodeObligation} value per graph node;
  \item a chain of \texttt{NodeObligation.compile} definitions;
  \item a final theorem \(\Gamma\vdash[\PRAAx]A\);
  \item an optional native corollary using \texttt{exportNative}.
\end{enumerate}
The generated source, not the generator, is checked by Lean.  This keeps the generator outside the trusted
base while preserving convenience.

\subsection{Validation sequence}

A disciplined integration run should execute:
\begin{enumerate}
  \item \texttt{lake env lean} on the three generic modules;
  \item \texttt{\#print axioms} on all exported compiler theorems;
  \item unit examples for substitution from arities zero, one, and several;
  \item examples with universal and existential formulas to exercise binder substitution;
  \item a three-node graph and a graph with a loop node;
  \item the cubic-bisection example;
  \item native export of the final ordinary-iPRA theorem.
\end{enumerate}
Only after these pass should large domain packets be migrated.  The main risk is API alignment, not the
mathematical proof of the structural kernel.
